% =====================================================================
%  preamble.tex -- packages and ALL custom macros for the B23 solutions
%  (keep every macro here; answer files should only use these)
% =====================================================================
\usepackage[T1]{fontenc}
\usepackage[utf8]{inputenc}
\usepackage{lmodern}
\usepackage[a4paper,margin=2.3cm,headheight=15pt]{geometry}
\usepackage{amsmath,amssymb,amsthm,mathtools,mathrsfs}
\usepackage{enumitem}
\usepackage{xcolor}
\usepackage{tikz}
\usetikzlibrary{arrows.meta,positioning,calc,cd}
\usepackage{tikz-cd}
\usepackage{graphicx}
\usepackage{array,booktabs,longtable}
\usepackage[most]{tcolorbox}
\usepackage{fancyhdr}
\usepackage{pgffor}
\usepackage{needspace}
\usepackage[colorlinks=true,linkcolor=blue!50!black,urlcolor=blue!50!black]{hyperref}
\usepackage{bookmark}

\setlength{\parindent}{0pt}
\setlength{\parskip}{0.5em}
\setlist{itemsep=2pt,topsep=3pt}

% ---------------------------------------------------------------------
%  Common math notation (students' answers may use these)
% ---------------------------------------------------------------------
\newcommand{\R}{\mathbb{R}}
\newcommand{\N}{\mathbb{N}}
\newcommand{\Z}{\mathbb{Z}}
\newcommand{\Q}{\mathbb{Q}}
\newcommand{\C}{\mathbb{C}}
\newcommand{\T}{\mathcal{T}}
\newcommand{\B}{\mathscr{B}}
\newcommand{\Pow}{\mathcal{P}}
\DeclareMathOperator{\Int}{Int}
\DeclareMathOperator{\Bd}{Bd}
\DeclareMathOperator{\id}{id}
\newcommand{\cl}[1]{\overline{#1}}
\newcommand{\onto}{\twoheadrightarrow}

% theorem-like environments for use inside solutions
\theoremstyle{plain}
\newtheorem*{theorem}{Theorem}
\newtheorem*{lemma}{Lemma}
\newtheorem*{proposition}{Proposition}
\newtheorem*{corollary}{Corollary}
\newtheorem*{claim}{Claim}
\theoremstyle{definition}
\newtheorem*{definition}{Definition}
\newtheorem*{example}{Example}
\theoremstyle{remark}
\newtheorem*{remark}{Remark}

% ---------------------------------------------------------------------
%  Page headers: left = problem set, right = problem (+ solver name)
% ---------------------------------------------------------------------
\pagestyle{fancy}
\fancyhf{}
\fancyhead[L]{\small\leftmark}
\fancyhead[R]{\small\tasrighthead}
\fancyfoot[C]{\thepage}
\fancyfoot[R]{\scriptsize Topology Assignment Solutions -- B23}
\renewcommand{\headrulewidth}{0.4pt}
\renewcommand{\sectionmark}[1]{}
\renewcommand{\subsectionmark}[1]{}
% Right header = the problem (and solver) running at the TOP of the page, so the later pages
% of a long solution keep showing its solver. Every \problem ends with an empty mark, so if
% the page starts with a fresh problem the top mark is blank and the first mark is shown.
\ExplSyntaxOn
\cs_new:Npn \tasrighthead
  {
    \tl_if_blank:eTF { \TopMark{2e-right} }
      { \FirstMark{2e-right} }
      { \TopMark{2e-right} }
  }
\ExplSyntaxOff

% ---------------------------------------------------------------------
%  Boxes
% ---------------------------------------------------------------------
\definecolor{stmtbg}{RGB}{240,244,250}
\definecolor{stmtframe}{RGB}{70,100,150}
\definecolor{pendingbg}{RGB}{255,243,224}
\definecolor{pendingframe}{RGB}{230,130,20}
\definecolor{solvedgreen}{RGB}{20,130,60}

\newtcolorbox{statementbox}{enhanced,breakable,colback=stmtbg,
  colframe=stmtframe,boxrule=0.5pt,arc=2pt,left=6pt,right=6pt,top=4pt,bottom=4pt}

\newcommand{\pendingbox}{%
  \begin{tcolorbox}[colback=pendingbg,colframe=pendingframe,boxrule=0.6pt,arc=2pt]
    \centering\bfseries\color{pendingframe!80!black} Solution pending
  \end{tcolorbox}}

% handwriting that could not be read with certainty
\newcommand{\illegible}[1]{\colorbox{yellow!45}{\textup{[illegible: #1?]}}}

% ---------------------------------------------------------------------
%  Structure macros
% ---------------------------------------------------------------------
\makeatletter
\newcommand{\tas@twodigit}[1]{\ifnum#1<10 0\fi\number#1}
\def\tas@currentproblem{}

% solver names collected from the previous run (via .aux) for the TOC
\newcommand{\tas@addsolver}[2]{%
  \@ifundefined{tas@solvers@#1}%
    {\expandafter\gdef\csname tas@solvers@#1\endcsname{#2}}%
    {\expandafter\g@addto@macro\csname tas@solvers@#1\endcsname{, #2}}}
\AtBeginDocument{\immediate\write\@auxout{\string\providecommand\string\tas@addsolver[2]{}}}

% \problemset{<n>}{<optional subtitle>}
\newcommand{\problemset}[2]{%
  \clearpage
  \phantomsection
  \section*{Problem Set #1\ifx\relax#2\relax\else\ \normalfont\large(#2)\fi}%
  \addcontentsline{toc}{section}{Problem Set #1\ifx\relax#2\relax\else\ (#2)\fi}%
  \markboth{Problem Set #1}{}}

% \problem{<set>}{<question>}{<statement>}
%   prints the statement, then inputs answers/ps<set>/q<NN>.tex if it exists,
%   otherwise a "Solution pending" box.
\newcommand{\problem}[3]{%
  \def\tas@currentproblem{#1.#2}%
  \bigskip
  \needspace{8\baselineskip}% keep heading with its statement
  \phantomsection
  \subsection*{Problem #1.#2}%
  \addcontentsline{toc}{subsection}{Problem #1.#2%
    \@ifundefined{tas@solvers@#1.#2}{}%
      {\enspace{\protect\small\protect\itshape -- \csname tas@solvers@#1.#2\endcsname}}}%
  \markright{Problem #1.#2}%
  \begin{statementbox}#3\end{statementbox}%
  \IfFileExists{answers/ps#1/q\tas@twodigit{#2}.tex}%
    {\input{answers/ps#1/q\tas@twodigit{#2}.tex}}%
    {\pendingbox}%
  \par
  \markright{}}% end-of-problem mark (see \tasrighthead)

% \submittedby{<Full Name>} -- first line of every answer file
\newcommand{\submittedby}[1]{%
  \par\medskip
  \noindent{\large\bfseries Solution}\hfill
  {\itshape submitted by}\ \textbf{#1}\par
  \vspace{-0.4em}\noindent\rule{\linewidth}{0.3pt}\par
  \markright{Problem \tas@currentproblem\ -- #1}%
  \immediate\write\@auxout{\string\tas@addsolver{\tas@currentproblem}{#1}}}

% \contributor{<Full Name>}{<questions solved>} -- rows of the contributors table
% (rows are collected first, then printed by \contributorrows inside the table)
\newcommand{\tas@rows}{}
\newcommand{\contributor}[2]{\g@addto@macro\tas@rows{#1 & #2 \\}}
\newcommand{\contributorrows}{%
  \ifx\tas@rows\@empty
    \multicolumn{2}{@{}l}{\itshape No submissions yet.} \\
  \else\tas@rows\fi}

% ---------------------------------------------------------------------
%  Coverage table: \coveragerow{<set>}{<number of questions>}
% ---------------------------------------------------------------------
\newcounter{tas@solved}
\newcounter{tas@total}
\newcounter{tas@solvedall}
\newcommand{\tas@mark}[2]{%
  \IfFileExists{answers/ps#1/q\tas@twodigit{#2}.tex}%
    {\stepcounter{tas@solved}\stepcounter{tas@solvedall}%
     \textcolor{solvedgreen}{\textbf{#2}}}%
    {\textcolor{black!35}{#2}}}
\newcommand{\coveragerow}[2]{%
  \setcounter{tas@solved}{0}\addtocounter{tas@total}{#2}%
  PS#1 &
  \foreach \q in {1,...,#2}{\tas@mark{#1}{\q}\ }%
  & \arabic{tas@solved}/#2 \\}
\newcommand{\coveragetotal}{%
  \midrule
  \multicolumn{2}{r}{\textbf{Total solved}} &
  \textbf{\arabic{tas@solvedall}/\arabic{tas@total}} \\}
\makeatother
